\documentclass[10pt,a4paper]{amsart}
\usepackage[margin=19mm]{geometry}
\usepackage{amsmath,amssymb,mathtools}
\usepackage[scr=rsfs,cal=euler]{mathalfa}
\usepackage{xfrac}
\usepackage[unicode,hidelinks,bookmarks=false]{hyperref}
\newcommand{\ie}{i.e.\ }
\newcommand{\cf}{cf.\ }
\newcommand{\resp}{respectively\ }
\newcommand{\Subsection}{\subsection}
\providecommand{\nobreakdash}{\nobreak\hbox{-}}
\setlength{\emergencystretch}{2em}
\multlinegap0pt
\usepackage{bm}
\usepackage{mathtools}
\usepackage{xcolor}
\usepackage[shortlabels]{enumitem}
\newtheorem{proposition}{Proposition}
\newtheorem{lemma}{Lemma}
\newtheorem{theorem}{Theorem}
\usepackage{cleveref}
\crefname{proposition}{Proposition}{Propositions}
\crefname{lemma}{Lemma}{Lemmas}
\crefname{theorem}{Theorem}{Theorems}
\newcommand{\norm}[1]{\left\|#1\right\|}
\title{Generalized differentiability and second-order necessary optimality conditions for an elliptic optimal control problem with exponential nonlinearity and discrete measures}
\author{Vu Huu Nhu, Nguyen Hai Son, Phan Quang Sang, Tran Duy}
\date{Source: arXiv:2605.19339v1 (CC BY 4.0)}
\begin{document}
\maketitle

\noindent\emph{Source and licence.} Generalized differentiability and second-order necessary optimality conditions for an elliptic optimal control problem with exponential nonlinearity and discrete measures. Version arXiv:2605.19339v1 (CC BY 4.0), distributed by arXiv under the Creative Commons Attribution 4.0 International licence, http://creativecommons.org/licenses/by/4.0/, as recorded in the arXiv licence metadata for this exact version and shown on the abstract page https://arxiv.org/abs/2605.19339v1. Full licence notice: \texttt{licenses/arxiv-2605.19339-CC-BY-4.0.txt}.\par\smallskip

\noindent\textit{Editorial scope.} Counted unit: the article's theorem thm:1st-NOCs (source0.tex lines 1324--1354). The article's complete original proof of it is delivered with it (source0.tex lines 1355--1383, one \texttt{proof} environment of 29 lines). No mathematics is written, repaired or re-derived: the statement and the proof are the article's own lines, in the article's own notation, and no retained line is substituted or re-flowed.  Retained uncounted as the source-local context the proof needs: lines 184--193; lines 1023--1055; lines 1288--1303.  Nothing was omitted from any retained span, so the package carries no omission record.  No retained line contains a citation, so the wrapper carries no bibliography and no bibliography entry.  No retained line contains a figure environment, an image-inclusion command or drawing code, so no image is delivered and the package declares no asset dependency.  Editorial formatting only, in addition: an AMS \texttt{amsart} preamble that restates the article's own declarations and macros used by the retained lines, namely the environments \texttt{proposition}, \texttt{lemma}, \texttt{theorem} on separate counters, together with the standard packages those lines need.  The retained environments print in this document as Proposition 1, Lemma 1, Theorem 1 in order of appearance, whereas the article prints the same items with the numbering of its own full document.  The complete native source of the article as distributed in the arXiv e-print for this exact version is bundled unchanged as \texttt{sources/200910/source0.tex}.
\begin{equation}
	\label{eq:P}
	\tag{P}
	\left\{
	\begin{aligned}
		& \min\limits_{\bm{u} = (u_i) \in \ell^1} J(\bm{u}) := \frac{1 }{2} \norm{y_{\bm{u}} - y_d}_{L^2(\Omega)}^2 + \frac{\nu }{2} \sum_{i=1}^\infty u_i^2 \\
		& \text{subject to } \bm{\alpha} \leq \bm{u} \leq \bm{\beta},
	\end{aligned}
	\right.
\end{equation}

\begin{proposition}
	\label{prop:cost-func-1st-2nd-directional-der}
	Assume that $\bm{u} \in \mathcal{D}$ such that $e^{S(\bm{u})} \in L^r(\Omega)$ for some $r > 1$. Let $k > 1$ be an integer.
	Then the reduced cost functional $J$ is directionally differentiable at $\bm{u}$ in every direction $\bm{h} \in \ell^1_k$ and satisfies
	\begin{equation}
		\label{eq:cost-func-1st-directional-der}
		J'(\bm{u})\bm{h} = \int_\Omega (S(\bm{u}) - y_d)z_{\bm{u}; \bm{h}} dx + \nu \sum_{i=1}^k u_i h_i 
		 = \sum_{i=1}^k \varphi_{\bm{u}}(x_i) h_i + \nu \sum_{i=1}^k u_i h_i
	\end{equation}
	with $z_{\bm{u}; \bm{h}} := S'(\bm{u})(\bm{h})$ and $\varphi_{\bm{u}}$ being the unique solution  in $W^{1,s_*}_0(\Omega)$ 
	for some $s_* :=  s_*(r) >2$	to 
	\begin{equation}
		\label{eq:adjoint-state0}
		\left\{
		\begin{aligned}
			-\Delta \varphi_{\bm{u}} + e^{S(\bm{u})} \varphi_{\bm{u}}& = S(\bm{u}) - y_d && \text{in } \Omega \\
			\varphi_{\bm{u}} &= 0 && \text{on } \Gamma.
		\end{aligned}
		\right.
	\end{equation}
	Moreover, there holds that
	\begin{equation}
		\label{eq:cost-func-2nd-directional-der-limit}
		\frac{1}{\rho^2} \left[J(\bm{u} + \rho \bm{h}) - J(\bm{u} ) - \rho J'(\bm{u})\bm{h} - \frac{\rho^2}{2} J''(\bm{u})\bm{h}^2   \right] \to 0 \quad \text{as } \rho \to 0^+,
	\end{equation}
	where $J''(\bm{u})\bm{h}^2$ is defined by
	\begin{equation} 
		\label{eq:cost-func-2nd-directional-der-def}
		J''(\bm{u})\bm{h}^2 % & := \int_{\Omega} (S(\bm{u}) - y_d)w_{\bm{u}; \bm{\omega}, \bm{\psi}} + z_{\bm{u}; \bm{\omega}}z_{\bm{u}; \bm{\psi}} dx + \nu \sum_{i=1}^k h_i^2 \notag \\
		  := \int_{\Omega} \left[1 - e^{S(\bm{u})} \varphi_{\bm{u}} \right] z_{\bm{u}; \bm{h}}^2 dx 	 + \nu \sum_{i=1}^k h_i^2.
	\end{equation}
	%We refer to $J''(\bm{u})\bm{h}^2$ as the \emph{second-order generalized derivative} of $J$ at $\bm{u}$ in the direction $\bm{h}$.
\end{proposition}

\begin{lemma}
	\label{lem:direction-approximations}
	Assume that $\bm{h} = (h_i) \in \ell^1$. Then, 
	$\bm{h} \in \overline{\mathrm{cone}}(U_{ad} - \bm{\bar u})$ if and only if there holds
		\begin{equation} 
			\label{eq:tangent-direction-equivalence}
			h_i 
			\left\{
			\begin{aligned}
				& \geq 0 && \text{if } \bar u_i = \alpha_i,\\
				& \leq 0 && \text{if } \bar u_i = \beta_i
			\end{aligned}
			\right. \quad \text{for all } i \geq 1.
		\end{equation} 
	Moreover, if $\bm{h} \in \ell^1$ fulfills \eqref{eq:tangent-direction-equivalence}, then   $\bm{h}_k  \in \mathrm{cone}(U_{ad} - \bm{\bar u})$.
\end{lemma}

\begin{theorem}[First-order necessary optimality conditions]
	\label{thm:1st-NOCs}
	Assume that $\bm{\bar u}$ is a local minimizer to \eqref{eq:P}. Then, there exists an adjoint state $\bar{\varphi} \in W^{1,s_*}_0(\Omega)$ for some $s_*>2$, 
	which, along with ${\bm{\bar u}}$ and $\bar y := S(\bar{\bm{u}})$, fulfills the system
	\begin{subequations}
		\label{eq:1st-NOCs}
		\begin{align}
			& \left\{
			\begin{aligned}
				-\Delta \bar y + (e^{\bar y} - 1) &= f_0(x) + \sum_{i =1}^\infty \bar u_i \delta_{x_i} && \text{in } \Omega, \\
				\bar y &= 0 && \text{on } \Gamma,
			\end{aligned}
			\right.  \label{eq:1st-NOC-state} \\
			& \left\{
			\begin{aligned}
				-\Delta \bar \varphi + e^{\bar y} \bar\varphi &= \bar y - y_d && \text{in } \Omega, \\
				\bar \varphi &= 0 && \text{on } \Gamma,
			\end{aligned}
			\right. \label{eq:1st-NOC-adjoint-state} 
			\intertext{and}
			& \left\{
			\begin{aligned}
				\bar\varphi(x_i) + \nu \bar u_i & \geq 0 && \text{if } \bar u_i = \alpha_i < \beta_i, \\
				\bar\varphi(x_i) + \nu \bar u_i & = 0 && \text{if } \bar u_i \in (\alpha_i, \beta_i), \\
				\bar\varphi(x_i) + \nu \bar u_i & \leq 0 && \text{if } \bar u_i = \beta_i > \alpha_i
			\end{aligned}
			\right.
			\quad \text{for } i \geq 1. \label{eq:1st-NOC-variational} 
		\end{align}
	\end{subequations}
\end{theorem}

\begin{proof}
	Obviously, $\bm{\bar u}$ and $\bar y$ fulfill \eqref{eq:1st-NOC-state}.
	For any $\bm{u} \in U_{ad}$, we set $\bm{h} := \bm{u} - \bm{\bar u}$ and thus $\bm{h} \in U_{ad} - \bm{\bar u}$.
	% \subset \overline{\mathrm{cone}}(U_{ad} - \bm{\bar u})$. 
	% By \cref{lem:direction-approximations}, we deduce $\bm{h}_k \in {\mathrm{cone}}(U_{ad} - \bm{\bar u})$ for all $k \geq 1$. 
	Fixing now $k$ yields
	\[
		\bm{u}_\rho := \bm{\bar u} + \rho \bm{h}_k \in U_{ad} \quad \text{for all } 0 < \rho < \rho_0 
	\]
	for some constant $\rho_0$ depending only on $k$. 	From this and the optimality of $\bm{\bar u}$, there holds
	\[
		\frac{1}{\rho} \left[ J(\bm{u}_\rho) - J(\bm{\bar u})\right] \geq 0 \quad \text{for all } 0 < \rho < \rho_0.
	\]
	In view of the $k$-dimensional directional differentiability of $J$; see \cref{prop:cost-func-1st-2nd-directional-der}, there holds
	\[
		J'(\bm{\bar u}) \bm{h}_k \geq 0,
	\]
	or, equivalently, 
	\[
		\sum_{i=1}^k \bar\varphi(x_i) (u_i - \bar u_i) + \nu \sum_{i=1}^k \bar u_i (u_i - \bar u_i) \geq 0,
	\]
	where ${\bar{\varphi}}$ satisfies \eqref{eq:1st-NOC-adjoint-state} according to \cref{prop:cost-func-1st-2nd-directional-der}.
	By taking $k \to \infty$, we obtain
	\begin{equation*}
		%\label{eq:1st-NOC-primal}
		\sum_{i=1}^\infty [\bar\varphi(x_i) + \nu \bar u_i] (u_i - \bar u_i) \geq 0 \quad \text{for all } \bm{u} = (u_i) \in U_{ad}.
	\end{equation*}
	This, together with a straightforward argument, gives \eqref{eq:1st-NOC-variational}.
\end{proof}

\end{document}
